% L9a student companion deck: European option pricing with BSM.
\documentclass[aspectratio=169,11pt]{beamer}
\usepackage{vnslides}

\newcommand{\E}{\mathbb{E}}
\newcommand{\Q}{\mathbb{Q}}
\newcommand{\gy}{g_y}
\newcommand{\ST}{S_T}

\title{\texorpdfstring{{\vnheavy L9a}}{L9a}: European Option Pricing with Black--Scholes--Merton}
\subtitle{CHEME 5660 Cornell University}
\author{Copyright Jeffrey Varner 2026. All Rights Reserved}

\begin{document}

\vntitlepage

\begin{frame}{Disclaimer and Risks}
  \vspace{0.6em}
  \textbf{\ink{Training only.}} This content is for training and informational
  purposes only. The teaching team does not make, give, or endorse any offer or
  solicitation to buy or sell securities or derivative products, or provide
  investment or trading advice or strategy.

  \vspace{1.1em}
  \textbf{\ink{Risk.}} Trading involves risk and is generally not recommended for
  those with limited resources, experience, or a low risk tolerance. Past
  performance does not guarantee future success.

  \vspace{1.1em}
  \textbf{\ink{Responsibility.}} You are responsible for your investment decisions
  based on your financial circumstances, objectives, risk tolerance, and
  liquidity needs.
\end{frame}

\begin{frame}{Objectives for Today}
  L8b treated a quoted premium as an input. Today we compute that premium for a
  contract exercisable only at expiration.

  \vspace{0.5em}
  {\large\textbf{\ink{Today's concepts}}}
  \begin{itemize}\setlength{\itemsep}{0.4em}
    \item \hilite{Explain risk-neutral valuation}: separate a no-arbitrage
      pricing measure from a real-world return forecast.
    \item \hilite{Price European calls and puts}: evaluate the
      Black--Scholes--Merton formulas with consistent units and conventions.
    \item \hilite{Check the premium}: use put--call parity, nonnegativity,
      strike monotonicity, and Monte Carlo standard errors.
  \end{itemize}
\end{frame}

\begin{frame}{Lecture and Example Stay on the Same Problem}
  \textbf{\ink{Lecture:}} European exercise, risk-neutral valuation, the BSM call
  and put formulas, input conventions, put--call parity, and numerical checks.

  \vspace{1.0em}
  \textbf{\ink{Example:}} Price the same European call and put by BSM and by
  risk-neutral Monte Carlo. Report simulation uncertainty, verify parity, and
  sweep both premiums across strike.

  \vspace{1.0em}
  \textbf{\ink{Next:}} L9b changes the exercise right. The European price from
  today becomes the benchmark for measuring the value of American early-exercise
  flexibility.
\end{frame}

\begin{frame}{European Exercise Means One Terminal Payoff}
  A European contract can be exercised only at expiration $T$. With strike
  $K>0$, terminal share price $\ST\geq0$, and $x^+:=\max(x,0)$, the per-share
  payoffs are
  \[
    \boxed{h_c(\ST;K)=(\ST-K)^+,
    \qquad h_p(\ST;K)=(K-\ST)^+.}
  \]
  \begin{itemize}\setlength{\itemsep}{0.35em}
    \item \hilite{Premiums}: $C_0$ for the call and $P_0$ for the put, both in
      USD/share at time $0$.
    \item \hilite{Deliverable}: multiply by shares/contract only when converting
      a per-share value to account dollars.
    \item \hilite{Scope}: there is no early-exercise decision to optimize. L9b
      adds that decision for American contracts.
  \end{itemize}
\end{frame}

\begin{frame}{Risk Neutral Is a Pricing Measure, Not a Forecast}
  Let $\gy$ be the continuously compounded risk-free rate, $\delta$ a continuous
  dividend yield, and $\sigma>0$ annualized volatility. Under the BSM pricing
  measure $\Q$,
  \[
    \frac{dS_t}{S_t}=(\gy-\delta)\,dt+\sigma\,dW_t^{\Q}.
  \]
  A European payoff $h(\ST)$ has model value
  \[
    \boxed{V_0=e^{-\gy T}\E^{\Q}[h(\ST)].}
  \]
  \begin{itemize}\setlength{\itemsep}{0.3em}
    \item The drift $\gy-\delta$ enforces no-arbitrage pricing under the model.
    \item No estimate of the asset's historical or expected real-world return
      enters this valuation.
  \end{itemize}
\end{frame}

\begin{frame}{The Closed Form Comes with Assumptions}
  Classical BSM assumes:
  \begin{itemize}\setlength{\itemsep}{0.45em}
    \item frictionless continuous trading and no arbitrage;
    \item constant risk-free rate $\gy$, dividend yield $\delta$, and volatility
      $\sigma$ over the contract's life;
    \item continuous lognormal underlying-price paths; and
    \item the ability to trade the underlying and financing asset as required by
      the replicating hedge.
  \end{itemize}
  \vspace{0.4em}
  Real markets have jumps, discrete trading, volatility smiles, funding spreads,
  transaction costs, and constraints. A formula can be evaluated exactly while
  its model remains approximate.
\end{frame}

\begin{frame}{Black--Scholes--Merton Prices the Pair}
  Define the dimensionless quantities
  \[
    d_1=\frac{\ln(S_0/K)+(\gy-\delta+\tfrac12\sigma^2)T}{\sigma\sqrt T},
    \qquad d_2=d_1-\sigma\sqrt T.
  \]
  With $\Phi$ the standard-normal CDF, the European premiums are
  \[
    \boxed{C_0=S_0e^{-\delta T}\Phi(d_1)-Ke^{-\gy T}\Phi(d_2)}
  \]
  \[
    \boxed{P_0=Ke^{-\gy T}\Phi(-d_2)-S_0e^{-\delta T}\Phi(-d_1).}
  \]
  Both values are in USD/share. Set $\delta=0$ for a non-dividend-paying
  underlying.
\end{frame}

\begin{frame}{Input Conventions Are Part of the Formula}
  Before evaluating BSM, reconcile every input:
  \begin{itemize}\setlength{\itemsep}{0.45em}
    \item \hilite{$S_0$ and $K$}: the same currency per share;
    \item \hilite{$T$}: time to expiration in years, with an explicit day-count
      convention;
    \item \hilite{$\gy$}: a continuously compounded annual rate;
    \item \hilite{$\sigma$}: annualized log-return volatility; and
    \item \hilite{$\delta$}: a continuous proportional dividend yield.
  \end{itemize}
  A mathematically correct formula with inconsistent conventions produces a
  precisely computed wrong answer.
\end{frame}

\begin{frame}{Put--Call Parity Checks the Closed Forms}
  The European call and put satisfy
  \[
    \boxed{C_0-P_0=S_0e^{-\delta T}-Ke^{-\gy T}.}
  \]
  \begin{itemize}\setlength{\itemsep}{0.4em}
    \item \hilite{Economic meaning}: call minus put replicates the
      dividend-adjusted share minus a bond paying $K$ at expiration.
    \item \hilite{Implementation use}: compute the call and put directly, then
      require the parity residual to match the implementation's precision.
    \item \hilite{Numerical caution}: do not obtain a tiny premium by subtracting
      nearly equal rounded quantities; cancellation can produce a negative value.
  \end{itemize}
\end{frame}

\begin{frame}{Monte Carlo Is a Second Route to the Same Value}
  Draw $M$ independent endpoints under $\Q$, discount each payoff
  $Y_k=e^{-\gy T}h(S_T^{(k)})$, and estimate
  \[
    \widehat V_M=\frac1M\sum_{k=1}^{M}Y_k,
    \qquad
    \widehat{\mathrm{SE}}(\widehat V_M)=\frac{s_Y}{\sqrt M}.
  \]
  \begin{itemize}\setlength{\itemsep}{0.35em}
    \item Compare $\widehat V_M-V_0$ in estimated standard-error units.
    \item Report the random-seed policy, number of samples, and standard error.
    \item A nominal 95\% interval is not guaranteed to contain the closed form on
      every run; interval membership is a diagnostic, not an identity.
  \end{itemize}
\end{frame}

\begin{frame}{Strike Moves Calls and Puts in Opposite Directions}
  Holding $S_0$, $T$, $\gy$, $\delta$, and $\sigma$ fixed:
  \begin{itemize}\setlength{\itemsep}{0.55em}
    \item \hilite{Calls are nonincreasing in $K$}: a higher purchase price can
      only reduce the terminal call payoff state by state.
    \item \hilite{Puts are nondecreasing in $K$}: a higher sale price can only
      increase the terminal put payoff state by state.
    \item \hilite{Both premiums remain nonnegative}: they are prices of
      nonnegative terminal payoffs.
  \end{itemize}
  \vspace{0.4em}
  The example notebook verifies these properties across strikes from 20 to 100
  USD/share.
\end{frame}

\begin{frame}{Four Checks Before Trusting the Premium}
  \begin{enumerate}\setlength{\itemsep}{0.5em}
    \item $C_0\geq0$ and $P_0\geq0$.
    \item Put--call parity holds within the formula's rounding tolerance.
    \item Calls decrease and puts increase as strike rises with other inputs fixed.
    \item Monte Carlo discrepancies are ordinary when measured in simulation
      standard-error units.
  \end{enumerate}
  \vspace{0.4em}
  These checks establish internal consistency under BSM. They do not prove that
  constant volatility, continuous paths, or the other assumptions describe a
  particular market well.
\end{frame}

\begin{frame}{Summary}
  L9a prices European options end to end: terminal payoff, risk-neutral
  valuation, BSM closed form, and numerical validation.

  \vspace{0.4em}
  \begin{itemize}\setlength{\itemsep}{0.4em}
    \item \hilite{Risk neutral is a pricing measure}: no real-world return
      forecast enters the premium.
    \item \hilite{BSM prices European exercise}: inputs and compounding
      conventions are part of the model specification.
    \item \hilite{A premium needs checks}: use direct formulas, parity,
      monotonicity, and Monte Carlo uncertainty together.
  \end{itemize}

  \vspace{0.5em}
  Next, L9b adds American exercise and asks when that extra right makes the
  American premium exceed this European benchmark.
\end{frame}

\end{document}
